\scanpage{62}
\mspage{61}
\mspage{35}

\deleted{\unsure{et}} $S'\longrightarrow S$. Posons dans cet situation
\[
\psi(S,Z,Y,X)=\xi\eta\varphi(Z,Y),
\]
d'où $\psi$ est de nature locale et commute aux changements de base. Or si $(Z,X)$ est un couple lisse de codimension $c>1$, on peut, localement sur $X$, trouver un triple $(Z,Y,X)$ avec $0<\operatorname{codim}(Y,X)<c$. Tout revient donc à démontrer que la flèche $\psi(S,Z,Y,X)$ est indépendant de $Y$. En effet, cela démontrera l'unicité de \qquad pour codimension $c$, et l'existence résultera de descente.

Si $(Z,Y,X)$ est arbitraire tel que $\operatorname{codim}(Y,X)>1$, on peut, localement sur $X$, trouver un triple lisse $(Y,W,X)$ où la dodimension de $(W,X)$ est égal à 1. Par une chasse de diagramme de transitivité, que nous laissons au lecteur\unsure{\}|,} on se ramène donc à vérifier l'indépendance de $Y$ lorsque $(Y,X)$ est de codimension 1. Soient \deleted{\unsure{$(Z,Y_1,X)$ et $(Z,Y_0$}} $(Z,Y_0,X)$ et $(Z,Y_1,X)$ deux tels triples. Puisque le problème est local, nous pouvons supposer que $Y_0$ et $Y_1$ sont définis chacun par une équation dans $X$, disons $f_0=0$ resp. $f_1=0$.

Soit $\overline S=\uline{\operatorname{Spec}}O_S[t]$, $\overline X=X\times_S\overline S$, $\overline Z=Z\times_S\overline S$, et soit $\overline Y$ le sous-schéma de $\overline X$ défini par l'équation
\[
(f_1-f_0)t+f_0=0.
\]
Alors $\overline Y/\overline S$ est lisse en les points au dessus des
